% ===========================================================================
% Gesamtskript VWL -- Übungsaufgaben Lektion 1: Grundbegriffe & Preisbildung
% (Fragment, wird per \input eingebunden; Lösungen in loesungen-l1.tex)
% Gesamtpunktzahl: 104 Punkte
% ===========================================================================

% ---------------------------------------------------------------------------
\subsection*{Aufgabe~1.1 --- Wirtschaften und ökonomisches Prinzip (11 Punkte)}
\addcontentsline{toc}{subsection}{Aufgabe~1.1 --- Wirtschaften und ökonomisches Prinzip}

Sie sind Mitarbeiterin bzw. Mitarbeiter der \emph{Nordlicht Fahrradmanufaktur GmbH} in Kiel,
die hochwertige E-Lastenräder herstellt. Eine Kundin, Frau Jansen, möchte ihre beiden Kinder
künftig ohne Auto zur Kita bringen. Sie hat dafür 4.500~€ angespart und bestellt schließlich
ein Lastenrad im Onlineshop der Nordlicht GmbH.

In der wöchentlichen Teamsitzung fallen folgende Aussagen:
\begin{itemize}
  \item[(1)] Einkaufsleiterin: „Für das nächste Quartal brauchen wir 2.400 Aluminiumrahmen.
    Ich hole Angebote ein und bestelle beim günstigsten Lieferanten.``
  \item[(2)] Vertriebsleiter: „Für die Fahrradmesse in Friedrichshafen stehen uns
    18.000~€ zur Verfügung. Damit wollen wir so viele Händlerkontakte wie möglich gewinnen.``
  \item[(3)] Geschäftsführer: „Unser Ziel ist klar: Wir produzieren mit den geringstmöglichen
    Kosten die größtmögliche Zahl an Rädern.``
\end{itemize}

\begin{enumerate}[label=\alph*)]
  \item \textbf{Erläutern} Sie am Beispiel von Frau Jansen die Begriffe Bedürfnis, Bedarf
    und Nachfrage. \hfill (3~P.)
  \item \textbf{Ordnen} Sie die Aussagen (1) bis (3) jeweils dem Maximalprinzip, dem
    Minimalprinzip oder keinem ökonomischen Prinzip zu und \textbf{begründen} Sie Ihre
    Zuordnung. \hfill (6~P.)
  \item \textbf{Ordnen} Sie die folgenden Gegebenheiten der Nordlicht GmbH den
    volkswirtschaftlichen Produktionsfaktoren zu: das zwölfköpfige Montageteam, das
    Gewerbegrundstück im Kieler Hafen, der neue Lackierroboter, das patentierte Verfahren
    zur Rahmenverklebung. \hfill (2~P.)
\end{enumerate}

% ---------------------------------------------------------------------------
\subsection*{Aufgabe~1.2 --- Marktarten und Marktformen (12 Punkte)}
\addcontentsline{toc}{subsection}{Aufgabe~1.2 --- Marktarten und Marktformen}

Sie sind Mitarbeiterin bzw. Mitarbeiter der \emph{Hanse Agrarhandel GmbH} in Lüneburg. Das
Unternehmen handelt mit Saatgut, Dünger und Pflanzenschutzmitteln und vermarktet die Ernte
seiner landwirtschaftlichen Kunden.

\begin{enumerate}[label=\alph*)]
  \item \textbf{Ordnen} Sie die folgenden Vorgänge jeweils der zutreffenden Marktart zu:
    \begin{enumerate}[label=(\arabic*)]
      \item Die Hanse Agrarhandel GmbH schreibt zwei Stellen für Lagerfachkräfte aus.
      \item Sie mietet für die Erntesaison eine zusätzliche Lagerhalle an.
      \item Sie nimmt bei ihrer Hausbank einen Investitionskredit über 250.000~€ auf.
      \item Sie verkauft 300~Tonnen Speisekartoffeln an eine Supermarktkette.
    \end{enumerate}
    \hfill (4~P.)
  \item \textbf{Bestimmen} Sie für die folgenden Situationen die Marktform und
    \textbf{begründen} Sie Ihre Antwort:
    \begin{enumerate}[label=(\arabic*)]
      \item Rund 1.800 Landwirte der Region liefern ihre Zuckerrüben an die beiden
        Zuckerfabriken in Norddeutschland.
      \item Eine einzige Reederei betreibt die Fährverbindung zu einer Nordseeinsel und
        befördert jährlich mehrere Hunderttausend Fahrgäste.
      \item Pflanzenschutzmittel werden von wenigen internationalen Chemiekonzernen
        hergestellt und an Tausende Händler und Landwirte verkauft.
    \end{enumerate}
    \hfill (6~P.)
  \item Ein Kollege behauptet: „Der Lüneburger Wochenmarkt ist ein vollkommener Markt, denn
    dort gibt es viele Anbieter und viele Nachfrager.`` \textbf{Beurteilen} Sie diese
    Aussage. \hfill (2~P.)
\end{enumerate}

% ---------------------------------------------------------------------------
\subsection*{Aufgabe~1.3 --- Preiselastizität und Umsatzwirkung (16 Punkte)}
\addcontentsline{toc}{subsection}{Aufgabe~1.3 --- Preiselastizität und Umsatzwirkung}

Sie sind im Controlling der \emph{Bohnenwerk Kaffeerösterei GmbH} in Bremen tätig. Im
letzten Quartal wurden in einem Testmarkt die Preise zweier Produkte erhöht. Alle übrigen
Einflussgrößen blieben unverändert. Es liegen folgende Daten vor (Absatz in Packungen je
Quartal):

\begin{center}
\begin{tabular}{l c c c c}
\toprule
\rowcolor{bfwTabHead}
\textbf{Produkt (je 500\,g)} & \textbf{Preis alt} & \textbf{Preis neu} & \textbf{Absatz alt} & \textbf{Absatz neu}\\
\midrule
Espresso „Crema Oro`` & 14,00~€ & 15,40~€ & 8.000 & 6.800\\
Filterkaffee „Hausmischung`` & 6,00~€ & 6,60~€ & 20.000 & 19.200\\
\bottomrule
\end{tabular}
\end{center}

\begin{enumerate}[label=\alph*)]
  \item \textbf{Berechnen} Sie für beide Produkte die direkte Preiselastizität der
    Nachfrage. Der Rechenweg ist darzustellen. \hfill (4~P.)
  \item \textbf{Interpretieren} Sie beide Ergebnisse und \textbf{begründen} Sie, warum
    die Kunden auf die Preiserhöhungen unterschiedlich stark reagieren. \hfill (4~P.)
  \item \textbf{Berechnen} Sie für beide Produkte den Umsatz vor und nach der
    Preiserhöhung sowie die Umsatzveränderung in Euro und in Prozent. \hfill (4~P.)
  \item Der Vertriebsleiter schlägt vor, die Preiserhöhung bei beiden Produkten dauerhaft
    und bundesweit umzusetzen. \textbf{Beurteilen} Sie diesen Vorschlag. \hfill (4~P.)
\end{enumerate}

% ---------------------------------------------------------------------------
\subsection*{Aufgabe~1.4 --- Marktgleichgewicht aus Angebots- und Nachfragefunktion (17 Punkte)}
\addcontentsline{toc}{subsection}{Aufgabe~1.4 --- Marktgleichgewicht aus Angebots- und Nachfragefunktion}

Sie sind Mitarbeiterin bzw. Mitarbeiter der \emph{Weserholz Energie GmbH}, die Holzpellets
in 15-kg-Säcken vertreibt. Für den regionalen Pelletmarkt, der annähernd die Bedingungen
der vollständigen Konkurrenz erfüllt, hat die Marktforschung folgende Funktionen ermittelt
($p$ = Preis in € je Sack, $x$ = Menge in Säcken pro Tag):

\begin{center}
Nachfrage: \quad $x_N = 1.200 - 40\,p$ \qquad\qquad Angebot: \quad $x_A = 60\,p - 300$
\end{center}

\begin{enumerate}[label=\alph*)]
  \item \textbf{Berechnen} Sie den Gleichgewichtspreis, die Gleichgewichtsmenge und den
    Umsatz im Gleichgewicht. \hfill (4~P.)
  \item \textbf{Berechnen} Sie für einen Preis von 12~€ und für einen Preis von 20~€
    jeweils die angebotene und die nachgefragte Menge. \textbf{Bestimmen} Sie Art und
    Höhe des jeweiligen Überhangs und \textbf{erläutern} Sie die Marktsituation
    (Käufer- oder Verkäufermarkt) sowie die zu erwartende Preisreaktion. \hfill (6~P.)
  \item \textbf{Stellen} Sie Angebots- und Nachfragekurve in einem Preis-Mengen-Diagramm
    \textbf{dar}. Kennzeichnen Sie das Marktgleichgewicht sowie den Überhang bei einem
    Preis von 20~€. \hfill (4~P.)
  \item \textbf{Berechnen} Sie die Konsumentenrente im Gleichgewicht und
    \textbf{erläutern} Sie, was diese Größe aussagt. \hfill (3~P.)
\end{enumerate}

% ---------------------------------------------------------------------------
\subsection*{Aufgabe~1.5 --- Kurvenverschiebungen und Kreuzpreiselastizität (16 Punkte)}
\addcontentsline{toc}{subsection}{Aufgabe~1.5 --- Kurvenverschiebungen und Kreuzpreiselastizität}

Sie arbeiten im Vertrieb der \emph{Alpenquell Getränke GmbH}, eines mittelständischen
Herstellers von Limonaden. Auf dem Limonadenmarkt ereignen sich folgende Entwicklungen:
\begin{itemize}
  \item[(1)] Der Deutsche Wetterdienst sagt einen außergewöhnlich heißen und langen Sommer
    voraus.
  \item[(2)] Der Weltmarktpreis für Zucker steigt innerhalb weniger Monate um 40~\%.
  \item[(3)] Der Preis für Limonade sinkt.
  \item[(4)] In der Branche setzt sich eine neue, deutlich schnellere Abfülltechnik durch.
\end{itemize}

\begin{enumerate}[label=\alph*)]
  \item \textbf{Erläutern} Sie für jede Entwicklung (1) bis (4), ob sie zu einer Bewegung
    auf einer Kurve oder zu einer Verschiebung einer Kurve führt. Geben Sie jeweils an,
    welche Kurve betroffen ist und in welche Richtung sie sich bewegt bzw. verschiebt.
    \hfill (8~P.)
  \item Die Entwicklungen (1) und (2) treten gleichzeitig ein. \textbf{Erläutern} Sie die
    Auswirkungen auf den Gleichgewichtspreis und die Gleichgewichtsmenge. \hfill (4~P.)
  \item Ein Wettbewerber senkt den Preis seines Eistees um 20~\%. Der Absatz der
    Alpenquell-Limonade geht daraufhin von 12.500 auf 11.750 Kisten pro Monat zurück.
    \textbf{Berechnen} Sie die Kreuzpreiselastizität und \textbf{bestimmen} Sie die
    Güterbeziehung zwischen Eistee und Limonade. \hfill (4~P.)
\end{enumerate}

% ---------------------------------------------------------------------------
\subsection*{Aufgabe~1.6 --- Funktionen des Marktpreises (10 Punkte)}
\addcontentsline{toc}{subsection}{Aufgabe~1.6 --- Funktionen des Marktpreises}

Sie sind im Einkauf der \emph{Holzbau Kramer GmbH} tätig, eines Zimmereibetriebs in
Thüringen. Durch einen Bauboom und eine stark gestiegene Exportnachfrage aus Übersee ist
die Nachfrage nach Konstruktionsvollholz sprunghaft gestiegen. Der Marktpreis ist binnen
eines Jahres von 420~€ auf 680~€ je Kubikmeter gestiegen. Die Sägewerke fahren
Sonderschichten, ein Investor plant ein neues Sägewerk in der Region. Mehrere kleine
Zimmereien, die die höheren Einkaufspreise nicht an ihre Kunden weitergeben können,
haben ihren Betrieb aufgegeben.

\begin{enumerate}[label=\alph*)]
  \item \textbf{Erläutern} Sie drei Funktionen des Marktpreises jeweils anhand der
    beschriebenen Situation. \hfill (6~P.)
  \item \textbf{Beschreiben} Sie die Marktsituation, die unmittelbar nach dem
    Nachfrageanstieg beim alten Preis von 420~€ je Kubikmeter bestand, und
    \textbf{erläutern} Sie, wie der Preismechanismus den Markt zu einem neuen
    Gleichgewicht führt. \hfill (4~P.)
\end{enumerate}

% ---------------------------------------------------------------------------
\subsection*{Aufgabe~1.7 --- Monopolistischer Absatzbereich im unvollkommenen Polypol (10 Punkte)}
\addcontentsline{toc}{subsection}{Aufgabe~1.7 --- Monopolistischer Absatzbereich im unvollkommenen Polypol}

Sie sind im Produktmanagement der \emph{Senfmanufaktur Brinkhaus GmbH} tätig. Auf dem
deutschen Senfmarkt konkurrieren viele Hersteller mit zahlreichen Marken um viele
Verbraucher. Die Brinkhaus GmbH verkauft ihren „Altmühltaler Grobkornsenf`` in einem
Steingut-Glas mit traditionellem Etikett für 2,60~€ je Glas. Eine Marktforschungsstudie
hat ergeben, dass die Stammkundschaft bei Preisen zwischen 2,20~€ und 3,10~€ je Glas
weitgehend treu bleibt.

\begin{enumerate}[label=\alph*)]
  \item \textbf{Bestimmen} Sie die vorliegende Marktform unter Berücksichtigung der
    Marktvollkommenheit und \textbf{erklären} Sie, wodurch die Brinkhaus GmbH einen
    preispolitischen Spielraum besitzt. \hfill (4~P.)
  \item \textbf{Beurteilen} Sie die Auswirkungen auf den Absatz, wenn die Brinkhaus GmbH
    den Preis (1) auf 2,95~€ bzw. (2) auf 3,40~€ je Glas anhebt. \hfill (4~P.)
  \item \textbf{Nennen} Sie zwei Maßnahmen, mit denen die Brinkhaus GmbH ihren
    preispolitischen Spielraum vergrößern kann. \hfill (2~P.)
\end{enumerate}

% ---------------------------------------------------------------------------
\subsection*{Aufgabe~1.8 --- Preisbildung im Oligopol und im Monopol (12 Punkte)}
\addcontentsline{toc}{subsection}{Aufgabe~1.8 --- Preisbildung im Oligopol und im Monopol}

Sie sind Mitarbeiterin bzw. Mitarbeiter der \emph{Baustoffhandel Lindner GmbH} in
Niederbayern. Zement wird in Ihrer Region von nur drei Zementwerken angeboten; Abnehmer
sind mehrere Hundert Bauunternehmen, Betonwerke und Baustoffhändler. Wegen der hohen
Transportkosten lohnt sich der Bezug aus weiter entfernten Werken kaum.

\begin{enumerate}[label=\alph*)]
  \item \textbf{Bestimmen} Sie die Marktform und \textbf{begründen} Sie Ihre Antwort.
    \hfill (2~P.)
  \item \textbf{Erläutern} Sie zwei für diese Marktform typische Verhaltensweisen der
    Anbieter in der Preispolitik. \hfill (4~P.)
  \item Eines der Werke erwägt, seinen Zementpreis um 10~\% zu senken, um Marktanteile
    zu gewinnen. \textbf{Beurteilen} Sie die Erfolgsaussichten dieser Strategie.
    \hfill (3~P.)
  \item Angenommen, eines der Werke übernimmt die beiden anderen. \textbf{Erläutern} Sie,
    wie sich dadurch seine Preissetzungsmöglichkeiten verändern, und \textbf{nennen} Sie
    zwei Grenzen seiner Preissetzungsmacht. \hfill (3~P.)
\end{enumerate}
